% Nouveau EGA V (cote 25): archive cote 25, site page 58 (PDF page 59).
% typescript. Draft compared: none.
% AI-transcribed by Claude Opus 5.5 (claude-opus-5-5, adaptive thinking, effort high), Message Batches API, batch ega25b (msgbatch_01A7n5YuXubUPNK1zsepi8RG), 2026-11-13; not yet checked by eye.
\scanpage{25/58}
\mspage{53}

10. \uline{Pinceaux de sections hyperplanes}, \uline{et fibrations de variétés éclatées}

10. 1. Soit Z l'ensemble P-exceptionnel dans $\overset{\hand{\vee}}{P}$ relativement à une
certaine propriété constructible\var{\inserted{P}}{P}{insertion}{ega25_25_058-v0001}, de sorte que Z est une partie construc%
tible de $\overset{\hand{\vee}}{P}$ . \hand{(}\var{\inserted[typed]{Supposons S=Spec(k),}}{Supposons S=Spec(k),}{insertion}{ega25_25_058-v0002}\hand{)} On sait (cf nº\hand{\unsure{11}}\var{\deleted{suivant}}{}{deletion}{ega25_25_058-v0003}, où nous \unsure{faisans|faisons} un ratrappage
% NOTE: a long diagonal pencil stroke in the left margin beside this and the next two lines
\var{\inserted{de choses}}{de choses}{insertion}{ega25_25_058-v0004} qui aurai\var{\over{t}{ent}}{ent}{writing over}{ega25_25_058-v0005} sans doute dû venir dans les nºs antérieurs) que pour que
l'on ait codim\var{\xout{\illegible{1}}}{}{typewriter strike-out}{ega25_25_058-v0006}(Z,$\overset{\hand{\vee}}{P}$ ) $\geqslant$ 2, il f et s que \hand{``}\var{\xout{t}}{}{typewriter strike-out}{ega25_25_058-v0007}toute droite assez
générale" L dans $\overset{\hand{\vee}}{P}$ ne rencontre pas Z, \var{\xout{xxxxx}}{}{typewriter strike-out}{ega25_25_058-v0008} ou encore $\overline{Z}$ , - et il suffit
qu'il existe \uline{une} droite L dans $\overset{\hand{\vee}}{P}$ ne rencontrant pas k. Lorsque \var{\xout{xxxxxxxx}}{}{typewriter strike-out}{ega25_25_058-v0009}
k est infini, il f donc et s qu'il existe une doite L dans $\overset{\hand{\vee}}{P}$ ne rencon-%
trant pas $\overline{Z}$ . On appelle "pinceau linéaire\var{\deleted{"}}{}{deletion}{ega25_25_058-v0010}\hand{(}\var{\inserted[typed]{de sections hyperplanes de X\hand{''}}}{de sections hyperplanes de X\hand{''}}{insertion}{ega25_25_058-v0011}\hand{)} défini par la\var{\xout{d}}{}{typewriter strike-out}{ega25_25_058-v0012}droite
L de $\overset{\hand{\vee}}{P}$\hand{,} le L-pré\unsure{s}chéma $\underline{Y}_L$ (définition valable pour un S quel\var{\xout{f}}{}{typewriter strike-out}{ega25_25_058-v0013}conque).
Ainsi, les réflexions précédentes, jointes aux résultats des nºs 8 et 9,
nous donnent des critères d'existence de tels pinceaux, ayant des fibres
$\underline{Y}_{\hand{\xi}}$ ( \hand{$\xi \in$}L) satisfaisant toutes la propriété P, tout d'abord dans le
cas \var{\deleted{deux}}{}{deletion}{ega25_25_058-v0014} \hand{(}\var{\inserted[typed]{où S est un}}{où S est un}{insertion}{ega25_25_058-v0015}\hand{)} corps de base infini. \var{\xout{Maisxprécédantxxxxxx}}{}{typewriter strike-out}{ega25_25_058-v0016} Compte tenu de
8.2.\hand{,} lorsque \var{\xout{X}}{}{typewriter strike-out}{ega25_25_058-v0017} pour tout cycle premier associé à X, on a dim f(\unsure{F})$>$0,
on peut (en prenant la propriété P' = P + condition de régularité
pour $\emptyset_{\hand{\xi}}$ ) exiger que le pinceau $\underline{Y}_L$ soit plat sur L. Dans le cas
où S est quelconque, on peut encore\var{\xout{x}}{}{typewriter strike-out}{ega25_25_058-v0018} souvent, par le procédé de 7.1. ,
construire un tel pinceau au dessus d'un voisinage ouvert d'un point
donné s de S, pourvu que k(s) soit infini et qu'on sache que Z est fermé
(ce qui sera assuré dans divers cas par les résultats du par.5 et
l'hypothèse de propreté pour X $\rightarrow$ S).

\var{\margin{|}}{\placedmargin{|}}{marginal note placed}{ega25_25_058-v0019} Pour bien faire, il conviendrait, après des explications géné-%
rales de cette eau, de donner de\var{\xout{x}}{}{typewriter strike-out}{ega25_25_058-v0020}s énoncés récapitulatifs où on
applique effectivement les résultats précédents pour un certain nombre
de propriétés de cette nature (y compris bien sur des propriétés modulées\hand{)}

\mspage{58}
